\section*{Validation}
\label{sec:validation}

Validation is the part of this literature most vulnerable to circularity,
and it is where the present analysis departs most sharply from earlier
versions of this work.  We first document a failure mode that inflates
reported accuracy by an order of magnitude, then give leakage-free
estimates of what the method can actually do.

\subsection*{Prior leakage in holdout designs}

A held-out text must contribute nothing to the model that dates it.  In
practice this is usually taken to mean exclusion from training and from
feature selection.  It must also mean exclusion from the prior, and this is
the requirement that is easy to violate without noticing.

Consider a text $u$ with scholarly date $d_u$ and stated uncertainty
$\sigma_u$.  If the posterior for $u$ is computed as
$p(d \mid \mathbf{x}_u) \propto \mathcal{L}(\mathbf{x}_u \mid d)\,
\mathcal{N}(d; d_u, \sigma_u^2)$, then the answer has been supplied as the
prior.  Under a Gaussian likelihood of precision $\tau_L$, the posterior
mean is a precision-weighted average in which the data receives weight
$\tau_L / (\tau_L + \sigma_u^{-2})$.  For a securely anchored text ---
small $\sigma_u$ --- that weight goes to zero, and the model returns
$d_u$ almost exactly.  The design therefore appears \emph{most} successful
precisely on the texts whose dates are best established, which is the
opposite of a stringent test.

This is not hypothetical.  Table~\ref{tab:leakage} gives the decomposition
for the four held-out units of the present corpus.  Backing the likelihood
precision out of the posterior and prior variances via
$\tau_L = \sigma_{\text{post}}^{-2} - \sigma_u^{-2}$ yields an implied
likelihood standard deviation of 48.6~yr (IQR 43.3--53.1) that is
remarkably stable across texts of very different length.  The weight
actually placed on the linguistic evidence, however, ranges from 2.8\% for
Haggai ($\sigma_u = 10$~yr) to 93.1\% for Zechariah~9--14
($\sigma_u = 150$~yr).  Across all 25 dated units the median ratio of
posterior to prior standard deviation is 0.851, and the median absolute
displacement of the MAP from the prior mean is 10.4~yr.

The consequence is severe.  Under this design the four holdouts return a
mean absolute error of 8.1~yr.  Dating the same texts from the likelihood
alone --- that is, from the linguistic evidence --- gives 103.7~yr.  We
report the latter.  We note that an independent implementation of the same
design for the Ancient Greek corpus reproduces the artifact exactly, with a
reported holdout MAE of 8.7~yr against 120.7~yr from the likelihood alone;
we therefore withdraw the cross-language validation claim rather than
repair it, since the corrected Greek result no longer supports the
generalisation it was offered as evidence for.

\subsection*{Leave-one-out design}

All results below use the following protocol.  For each of the 25 dated
units in turn, the unit is removed; feature screening (Spearman, $\alpha =
0.05$), standardisation, and all model fitting are performed on the
remaining 24; and the held-out unit is then dated under an agnostic prior.
No quantity derived from a unit influences its own prediction.  Two model
families are fitted and both are reported: the generative family of
Eq.~\ref{eq:loglik}, which inverts per-feature regressions, and a ridge
regression of date on screened features with the penalty chosen by inner
leave-one-out.  Significance is assessed by permuting the date vector and
re-running the entire pipeline, 2000 draws, so that the null distribution
absorbs selection effects exactly as the observed analysis does.

\subsection*{What the signal supports}

The diachronic signal is real.  Pairwise ordering accuracy --- given two
texts, does the model place them in the correct chronological order? --- is
69.8\% of 295 pairs for the generative family and 67.8\% for ridge, against
a 50\% baseline ($p = 4 \times 10^{-12}$ and $5 \times 10^{-10}$,
binomial).  Leave-one-out Spearman correlation between true and predicted
date is $+0.575$ (generative) and $+0.494$ (ridge), against permutation
nulls centred near zero ($p = 0.016$, $0.024$).  This is the one result
that replicates across both families at the level of the primary endpoint,
and it is the strongest finding in the paper.

\subsection*{What the signal does not support}

Absolute dating does not survive contact with an honest baseline.
Leave-one-out MAE is 156.2~yr for the generative family and 120.9~yr for
ridge.  Assigning every text the corpus mean date of 503~BCE gives 137.3~yr.
The generative family --- the one used for the point estimates reported in
earlier versions of this work --- is therefore \emph{worse} than a
constant, and ridge improves on a constant by 16~yr.  Any MAE figure in
this literature should be read against roughly 137~yr, not against zero.

Period assignment is correspondingly weak.  Four-way exact accuracy is
40.0\% (generative) and 48.0\% (ridge) against a 36\% majority-class
baseline; pre-exilic versus post-exilic classification reaches 72.0\% and
76.0\% against a 68\% baseline, significant in neither family
($p = 0.107$, $0.070$).  Across ten tests --- five metrics in two families,
strongly correlated --- nothing survives Benjamini-Hochberg correction at
$q = 0.05$, and four results survive at $q = 0.10$.  We report this
directly: no individual accuracy claim in this paper is robust to
multiplicity correction, and the ordinal result is defended on replication
across model families rather than on any single $p$-value.

The failure is structured rather than random.  The Exilic period is a
47-year window set against roughly 130~yr of irreducible error, and it is
simply not resolvable: of four exilic texts, the generative family places
none correctly and ridge one.  Hellenistic units collapse into Persian.
Pre-exilic is the only period the models recover reliably, and it is
recovered because it is wide and terminal, not because early Hebrew is
distinctive in some way the interior boundaries are not.

\subsection*{Calibration}

The parametric intervals used in earlier versions of this work achieve
52\% empirical coverage at a nominal 68\% (61\% under the residual
formulation; the manuscript previously reported 32\% for the original
specification).  The mechanism is visible in the multiverse analysis:
chance-correlated features receive large fitted coefficients and small
residual variances, and because the multivariate normal likelihood weights
each feature by $1/\sigma_j$, adding noise features \emph{narrows} the
credible interval.  Adding six and then twelve pure-noise features to a
clean specification shrinks the 68\% HDI for the P source from 329 to 285
to 246~yr while moving the MAP by less than 50~yr.  Interval width in this
framework is not a reliable index of evidence.

We therefore replace parametric intervals with conformal prediction
intervals calibrated on the leave-one-out residuals of the dated corpus.
These are distribution-free and carry guaranteed finite-sample marginal
coverage under exchangeability.  The resulting half-widths are $\pm233$~yr
(generative) and $\pm174$~yr (ridge) at 68\%.  These are the intervals
reported in Table~\ref{tab:targets}.  Exchangeability between the
calibration units and the targets is an assumption and not a guarantee:
the calibration units are whole books or coherent sections, whereas several
targets are cross-book composites.  Residual magnitude shows no detectable
dependence on text length ($\rho = -0.23$, $p = 0.27$), which supports the
use of a single conformal width.
